%% Complete new localization proofs; released source01 annex remains frozen.

\begin{annexitem}{annex:p10.lawvere-localization.denominators}{thm:p10-lawvere-localization-denominators}{The exact u-star/h filter and both denominator classes}
\begin{proof}
Finite-product preservation transports every finitary
internal Heyting operation and its equations on \(T\)
to the declared algebra \(u_*(T)\).  Composing global
predicates with the declared homomorphism \(h\) therefore
preserves bottom, top, meet, join and implication.  Thus
\(\theta_h\) is the stated homomorphism of global truth
algebras.  Its inverse image of top contains top and is
closed under meets; if \(U\leq V\) and \(\theta_h(U)\)
is top, their meet equation forces \(\theta_h(V)\) to
be top too.  This proves the filter conditions.
If bottom belongs to it, preservation of bottom gives
\(\mathsf{false}_0=\mathsf{true}_0\), which identifies
the initial and terminal subobjects of \(1_0\) and
makes \(\mathcal E_0\) degenerate.  Otherwise it is
proper.  No primality, finiteness or external choice
was used.

For a mono \(m\), the already proved adjunction on
monos identifies \(\Pi_{X\to1}(m)\) with
\(\forall_{X\to1}(m)\), as a subobject of \(1\).
Support of this mono is itself.  This proves exact
agreement with the source condition, without assuming
that \(u_*\) preserves quantifiers or classifiers.

For \(U\in\mathcal F_h\), the predicate adjunction
gives \(U\leq\forall_{X\to1}(U\times X)\).
Upward closure therefore puts \(i_{U,X}\) in the
source class.  Conversely the right-adjoint counit
gives \(\gamma_m\times X\leq A\) over \(X\).
Pulling \(m\) back to \(\gamma_m\times X\) is
the identity subobject.  Hence
\(\gamma_m\times A\to\gamma_m\times X\) is an
isomorphism.  In any category inverting the uniform
restrictions, both projections to \(A\) and \(X\)
are inverted, so their commuting square also inverts
\(m\).  The converse is immediate because all uniform
restrictions are source denominators.  This proves
equivalence of the two inversion requirements.

The source mono class contains identities, is closed
under composition and pullback: uniform admitted
subterminals witnessing two denominators can be
intersected, and their product restriction pulls
back to the same subterminal on any new domain.
The counit inclusion just proved gives a cofinal
uniform restriction of every denominator span.
These are exactly the properties needed for the
explicit fraction construction in the next item.
\end{proof}
\end{annexitem}

\begin{annexitem}{annex:p10.lawvere-localization.fractions}{thm:p10-lawvere-localization-fractions}{Fractions, all coherence equations, and the universal functor}
\begin{proof}
\paragraph{Sets, equivalence and composition.}
The disjoint union of the hom sets
\(\mathcal E(U\times X,Y)\), indexed by the set
\(\mathcal F_h\subseteq\operatorname{Sub}(1)\),
is a set in the stated universe.  Its quotient
therefore constructs each hom set without any
infinite colimit in \(\mathcal E\).
The relation is reflexive and symmetric.  For
transitivity intersect the two admitted witnesses
and both representative denominators; the resulting
admitted subterminal restricts both equalities
to the same maps.  This proves transitivity.
Restriction maps compose exactly: maps between
subterminals over \(1\) are unique, and the
induced product maps have their uniquely specified
projections.

For composition, all representatives and all
different candidate intersection objects can be
restricted to one further admitted intersection.
On it the two proposed composites are the same
ordinary composite with the same forced parameter
maps.  This proves independence of witnesses and
representatives.  For three arrows use the common
intersection of their denominators; ordinary
associativity proves the two germ composites equal.
The identity denominator \(1\) acts by the usual
product identification, so both unit laws hold.
This constructs the category and the functor \(Q\).

The projection \(i_{U,X}\) has inverse germ
\([U,1_{U\times X}]:X\to U\times X\).
Their composites equal identity after restriction
to \(U\); here \(U\times U\cong U\) because
\(U\) is subterminal.  Thus \(Q\) inverts uniform
denominators and, by the preceding lemma, every
source denominator.
Every germ is exactly the fraction
\(Q(f)Q(i_{U,X})^{-1}\).  A general source span
\(X\leftarrow A\to Y\) restricts through
\(\gamma_m\times X\to A\) to such a uniform
span.  Pullbacks give the common right-fraction
refinements, and the displayed equality relation
is precisely equality after further refinement.
There is no additional abstract Ore condition
left unproved.

\paragraph{The full universal property.}
Let \(H\) invert the source class.  It inverts
all \(i_{U,X}\).  If \(W\leq U\), the map
\(j:W\times X\to U\times X\) is inverted
as well, because \(i_{U,X}j=i_{W,X}\).
Equality of representatives on \(W\) consequently
equates
\(H(f)H(i_{U,X})^{-1}\) and
\(H(g)H(i_{V,X})^{-1}\).
Thus the stated formula is well defined.

For a composite, let \(W\leq U\cap V\) and
let \(k:W\times X\to V\times Y\) be the
forced parameter map paired with \(f|W\).
Its projection satisfies \(i_{V,Y}k=f|W\).
After applying \(H\) and inverting that
projection, this says
\(H(k)=H(i_{V,Y})^{-1}H(f|W)\).
Substituting and using
\(i_{U,X}j_{W,U}=i_{W,X}\) proves the
composition formula for \(\bar H\).  Identity
preservation is immediate at denominator \(1\).
Hence \(\bar H\) is a functor extending \(H\).
Since every germ is the displayed fraction,
every such extension has this same value; it
is unique on the unchanged objects and all
arrows.

A natural transformation \(H\to K\) has
its usual components at each object.  Naturality
for the numerator and for an invertible denominator
proves naturality for their fraction, so those
components extend uniquely.  The same argument
handles identity and composition of transformations.
This gives the functor-category localization
universal property, not only an objectwise
construction.  The preceding denominator lemma
makes it the localization of the exact original
\(u_* / h\) class.

If zero is admitted, refine every pair of
representatives to zero.  Its products are
initial by cartesian closedness, so all arrows
between each pair of objects coincide.  Every
object is then isomorphic to every other,
giving the degenerate category; no nondegeneracy
assumption is hidden in the construction.
\end{proof}
\end{annexitem}

\begin{annexitem}{annex:p10.lawvere-localization.logical}{thm:p10-lawvere-localization-logical}{Every topos structure and the unique localized classifier}
\begin{proof}
\paragraph{One finite-diagram refinement principle.}
A finite collection of representative germs can
be restricted to the intersection of its admitted
denominators.  A finite collection of equalities
of germs has, by definition, a finite collection
of further admitted witnesses; intersect them
too.  All those equations then hold as actual
maps on that one subterminal.  This includes
commuting triangles and two inverse identities.
Maps with parameter \(U\) can be lifted to
maps over \(U\) by pairing with its unique
parameter coordinate.  This proves the principle
for arbitrary finite diagrams and isomorphisms,
without an infinite choice of representatives.

\paragraph{Finite limits.}
The original terminal object is terminal after
\(Q\): each representative into it is the
unique ordinary map, and representatives agree
on any common admitted restriction.  For products,
two germs from \(Z\) to \(X,Y\) have a common
admitted denominator.  Pair their representatives
there to obtain a germ into \(X\times Y\).
Equality of its two projections gives equality
of the germ after a further common restriction.
This proves existence and uniqueness, and
preservation of the original product cone.

For original parallel maps \(f,g:X\rightrightarrows Y\),
take their original equalizer \(E\hookrightarrow X\).
A germ \([U,k]:Z\to X\) equalized by \(Qf,Qg\)
has a further admitted \(W\leq U\) where
\(fk|W=gk|W\).  It factors uniquely through
\(E\) there.  Monicity and the same refinement
criterion give uniqueness of its germ factor.
Thus \(Q\) preserves original equalizers.
For genuinely localized parallel arrows, first
represent both as ordinary maps
\(F,G:U\times X\rightrightarrows Y\) on one
admitted \(U\).  Their ordinary equalizer,
followed by the inverted projection to \(X\),
gives the equalizer in the localization by the
case just proved.  Hence all finite limits
exist, not only limits of strict source diagrams.

\paragraph{Finite colimits.}
The original initial object remains initial:
\(U\times0=0\) since product is a left
adjoint, so each hom germ out of it is unique.
Product with \(U\) preserves coproducts.
Consequently germs from \(X+Y\) to \(Z\)
are exactly pairs of germs from \(X,Y\),
by taking a common denominator for the pair.
This proves coproduct universality and
preservation by \(Q\).

Let \(q:Y\to C\) be an ordinary coequalizer
of \(f,g:X\rightrightarrows Y\).  A germ
\([U,k]:Y\to Z\) coequalizing their images
can be refined to \(W\) where the equation
is ordinary.  Since \(W\times(-)\) is a
left adjoint, \(W\times q\) is their
ordinary restricted coequalizer.  The map
\(k|W\) therefore factors through it.
For uniqueness, equality after \(Qq\)
has a further admitted witness where the
two maps agree after \(W\times q\);
cancel that epi to equate them there.
Thus \(Q\) preserves coequalizers.
For two arbitrary parallel germs, represent
them on \(U\times X\) as in the equalizer
case.  This source object is isomorphic to
\(QX\); the preserved ordinary coequalizer
is their required coequalizer.  Initial
objects, binary coproducts and coequalizers
generate finite colimits, supplying the
whole original finite-colimit package.

\paragraph{Exponentials and actual evaluation.}
For each \(U\), ordinary transposition gives
\[
 \mathcal E(U\times Z,Y^X)
       \cong\mathcal E(U\times Z\times X,Y).
\]
These bijections commute with restriction
of \(U\), since transposition is natural.
They therefore descend to inverse germ
bijections
\(\operatorname{Hom}(Z,Q(Y^X))\cong
\operatorname{Hom}(Z\times QX,QY)\).
The evaluation is \(Q\) of the original
evaluation, and the two transpose inverse
identities hold after a common restriction.
This proves cartesian closedness and its
preservation, including empty objects and
the improper-filter case.

\paragraph{All monomorphisms are eventually strict.}
An ordinary mono remains monic by finite-limit
preservation.  Let a localized mono be
represented by \(m:U\times X\to Y\).
Take its ordinary kernel pair with projections
\(p,q\).  Their images are equal because
the image of \(m\) is monic.  The equality
of those germs has a witness \(V\in\mathcal F_h\).
Restricting further to \(V\cap U\), its
kernel-pair projections agree, so the
restricted map to \(V\times Y\) is monic.
Its composite to \(Y\) is still monic,
because \(V\) is subterminal.  Its domain
is isomorphic to \(QX\) in the localization.
Thus every localized mono is represented,
up to that domain isomorphism, by an actual
mono of \(\mathcal E\).  This supplies
the missing direction of the classifier
construction, not just preservation of
initially chosen mono representatives.

\paragraph{Classifier existence and uniqueness.}
For such an ordinary mono \(M\hookrightarrow Y\),
its original characteristic map \(\chi_M:Y\to T\)
and truth pullback are preserved by the finite
limit proof.  It therefore classifies the
localized mono.  Now suppose a germ
\([U,a]:Y\to T\) gives the same truth
pullback.  Its pullback is the image under
\(Q\) of the ordinary mono
\(N=a^*(\mathsf{true})\hookrightarrow U\times Y\),
followed by the inverted projection.  These
two mono domains are isomorphic over \(QY\).
By the finite-diagram refinement principle,
the isomorphism, its inverse, their two
inverse equations and their maps to \(Y\)
are all realized over some admitted
\(V\leq U\).  Thus the corresponding
ordinary subobjects of \(V\times Y\)
agree.  Original classifier uniqueness
equates \(a|V\) with \(\chi_M|V\).
Their germs are equal.  This proves the
unique characteristic-map property for
every localized mono.

We have now proved every original topos
structure and its preservation.  In particular
\(QT\) is the classifier and \(Q\) is
logical, without proper/prime/finite-filter
assumptions.  The improper case is the
degenerate topos, where these universal
properties still hold.

\paragraph{The full internal logical constructions.}
Power objects and membership are constructed
from classifier exponentials and evaluation,
so their actual universal maps are preserved.
The earlier general dependent-product proof
uses only those power objects, finite equalizers
and the characteristic maps of its monic
functional graph.  Each such construction
has just been preserved, so it identifies
\(Q\Pi_f(E)\) with \(\Pi_{Qf}(QE)\),
including counit, inverse transposition and
parameter pullback maps.
Universal predicate quantification is its
restriction to monos, hence is preserved.
Every original epi is a kernel-pair
coequalizer by the proved exactness result;
finite coequalizer preservation therefore
preserves epis.  Finite limits preserve
monos, so epi--mono images are preserved
by their uniqueness in the now-proved
localized topos.  This preserves existential
predicate quantification and its internal
power-object map.  Heyting operations are
the already supplied classifier/finite-limit/
image constructions, so they too are
preserved.  No illicit exchange of
\(u_*\) with quantifiers was used.
\end{proof}
\end{annexitem}
